% !TEX root = Thesis_template.tex
%% appendix_code_architecture.tex
%% ---------------------------------------------------------------------
%% Appendix B: architecture of the code of Campaign 9 and of its
%% post-analysis. Rewritten 28 September 2026; the architecture figure is
%% drawn in TikZ here, so it needs no external tool.
%% ---------------------------------------------------------------------

\chapter{Architecture of the simulation code}\label{app:code-architecture}

\providecommand{\file}[1]{\expandafter\nolinkurl\expandafter{\detokenize{#1}}}

\section{Purpose and scope}\label{sec:arch-scope}

This appendix documents the code that produced the results of
Campaign 9 and of its post-analysis (Sections~\ref{sec:res-c9-formulation}
to~\ref{sec:res-c9-cont-front}), the formulation on which
the conclusions are based. It is written for a reader who needs to locate the
script behind a result, or to repeat a step of the post-analysis. The
earlier campaigns used earlier versions of the same modules, which the
history of the repository keeps.

The code is public, at
\url{https://github.com/dcrlopes/master-thesis-unipi}. The version
described here is version 1.0, the tag \texttt{v1.0} of the repository,
commit \texttt{283071d} of 2 October 2026, which later work on the
repository does not change. It contains 213 Python scripts and 33
shell scripts, and only those that compute a result of Campaign 9 or of
its post-analysis are described. Four groups of scripts are omitted:

\begin{itemize}[nosep]
  \item the shell scripts that launch, queue and monitor the
        calculations,
  \item the patch scripts, \texttt{apply\_*.py} and \texttt{fix\_*.py},
        that modified the source between campaigns,
  \item the figure scripts, \texttt{make\_*.py},
        \file{c9_thesis_figures.py} and \file{c9_ceiling_two_pressures.py},
        which only draw stored results,
  \item the tests, the diagnostics and the scripts of Campaigns 1 to 8.
\end{itemize}

\section{Macro architecture}\label{sec:arch-macro}

Figure~\ref{fig:code-architecture} shows the blocks of the code and the
files that cross each interface:

\begin{itemize}[nosep]
  \item \textbf{Core model:} the modules that build the OpenMC
        models, imported by every script that executes a transport
        calculation.
  \item \textbf{Design evaluator:} the only place where a design
        vector becomes a transport result inside the loop.
  \item \textbf{Optimiser:} the surrogate-assisted loop, which does
        not call OpenMC, together with the scripts of the
        continuation of the search.
  \item \textbf{Post-analysis:} the scripts that re-evaluate the
        designs of the campaign outside the loop, split into those that
        execute OpenMC and those that only read stored results.
\end{itemize}

Purple blocks call OpenMC~\cite{romano_openmc_2015,romano_2021_depletion},
the coral block is the optimiser, and grey blocks are data, analysis
and outputs. An unlabelled arrow between two blocks of scripts is a
Python import.

\begin{figure}[htbp]
  \centering
  \definecolor{archPhysFill}{HTML}{EEEDFE}\definecolor{archPhysLine}{HTML}{534AB7}
  \definecolor{archOptFill}{HTML}{FAECE7}\definecolor{archOptLine}{HTML}{993C1D}
  \definecolor{archInfFill}{HTML}{F1EFE8}\definecolor{archInfLine}{HTML}{5F5E5A}
  \begin{tikzpicture}[
      font=\footnotesize,
      blk/.style={draw, rounded corners=3pt, line width=0.6pt, align=left,
                  inner sep=4pt, text width=#1, anchor=north},
      blk/.default=3.8cm,
      phys/.style={fill=archPhysFill, draw=archPhysLine},
      opt/.style={fill=archOptFill, draw=archOptLine},
      inf/.style={fill=archInfFill, draw=archInfLine},
      lab/.style={font=\scriptsize, inner sep=1.5pt, align=center},
      arr/.style={-{Stealth[length=5pt]}, line width=0.6pt, draw=archInfLine}]
    % columns: left x = 1.55, middle x = 5.6, right x = 11.0
    \node[blk=5.4cm, inf] (env) at (5.6, 0) {\textbf{Environment and nuclear data}\\
      OpenMC 0.15.3 in \texttt{openmc-env}, ENDF/B-VII.1};
    \node[blk, phys, below=0.75cm of env] (model) {\textbf{Core model}\\\ttfamily\scriptsize
      reactor\_model.py\\ zoning.py\\ core\_geometry.py\\ leu\_policy.py\\ hardware3d.py};
    \node[blk, phys, below=0.9cm of model] (eval) {\textbf{Design evaluator}\\\ttfamily\scriptsize
      openmc\_evaluator.py\\ boron\_objective.py};
    \node[blk=3.15cm, inf, anchor=center] (kt) at (1.45,0 |- eval.center) {\textbf{Reactivity target}\\\ttfamily\scriptsize
      ktarget\_table\_c8.json\\ \rmfamily from Campaign 8};
    \node[blk, opt, below=1.3cm of eval] (opt) {\textbf{Optimiser}\\\ttfamily\scriptsize
      reactor\_optimization.py\\ run\_optimization.py\\[3pt]
      \rmfamily\footnotesize Continuation of the search\\\ttfamily\scriptsize
      axial\_ratio\_model.py\\ seed\_c9\_axial.py\\ seed\_c9\_floor.py};
    \node[blk=4.1cm, phys] (ptr) at (11.0,0 |- model.north) {\textbf{Post-analysis with transport}\\\ttfamily\scriptsize
      dep\_common.py\\ c9\_dep\_replicas.py\\ c9\_dep\_core2d.py\\ c9\_dep\_core3d.py\\
      axial\_shape\_c9.py\\ confirm3d.py\\ ladder\_core\_c9.py\\ mtc\_scan.py\\
      boron\_worth.py\\ c9\_eval\_predicted\_front.py};
    \node[blk=4.1cm, inf, below=0.9cm of ptr] (pnt) {\textbf{Post-analysis without transport}\\\ttfamily\scriptsize
      c9\_front.py\\ c9\_step0.py\\ c9\_surrogate\_cv.py\\ c9\_gp\_predicted\_fronts.py\\
      c9\_acquisition\_replay.py\\ c9\_replay\_fidelity.py\\ c9\_reverse\_retro.py\\
      c9\_peaking\_2d3d.py\\ c9\_axial\_checks.py\\ c9\_axial\_correction.py\\
      c9\_axial\_rescore.py\\ mtc\_front\_table.py\\ c9\_repro\_check.py};
    \node[blk, inf, anchor=south] (thesis) at (opt.south |- pnt.south) {\textbf{Dissertation}\\
      figures and tables of Chapter~\ref{ch:results}};
    % imports and data
    \draw[arr] (env) -- node[lab, right=2pt] {cross sections, chain} (model);
    \draw[arr] (model) -- (eval);
    \draw[arr] (model.east) -- (model.east -| ptr.west);
    \draw[arr] (kt) -- (eval);
    \draw[arr, {Stealth[length=5pt]}-{Stealth[length=5pt]}] (eval) --
      node[lab, left=2pt] {$\mathbf{x}$\\$(F_{\Delta H}, c_\mathrm{max}, \mathbf{g})$} (opt);
    % continuation: eight-layer cycle lengths back to the optimiser
    \coordinate (cA) at ([yshift=-2.9cm]ptr.north west);
    \coordinate (cB) at ([yshift=-0.5cm]opt.north east);
    \draw[arr, dashed] (cA) -- (cA -| 7.95,0) -- (7.95,0 |- cB) -- (cB);
    \node[lab, anchor=east] at ([xshift=-0.05cm, yshift=0.95cm]cB -| 7.95,0) {eight-layer\\cycle lengths};
    % checkpoints to both post-analysis blocks, through one bus
    \coordinate (bA) at ([yshift=-0.9cm]opt.east);
    \draw[arr] (bA) -- (bA -| pnt.west);
    \draw[arr] (bA -| 8.5,0) -- node[lab, sloped, above] {checkpoints}
      (8.5,0 |- ptr.south) -- ([yshift=0.4cm]ptr.south west);
    % results
    \draw[arr] (ptr) -- node[lab, right=2pt] {\texttt{runs.json}\\\texttt{summary.json}} (pnt);
    \draw[arr] (thesis.east -| pnt.west) -- node[lab, above] {tables\\figures} (thesis.east);
  \end{tikzpicture}
  \caption{Code of Campaign 9 and of its post-analysis: macro architecture.}
  \label{fig:code-architecture}
\end{figure}

Campaign 9 and its post-analysis traverse the blocks in this order:

\begin{enumerate}[nosep]
  \item The optimiser samples the design of experiments by Latin-hypercube
        sampling~\cite{mckay_lhs_1979} and sends each design vector to
        the evaluator.
  \item The evaluator builds the OpenMC models, executes the assembly
        depletion and the core solves, and returns the objectives and
        the constraint vector.
  \item The optimiser fits the Gaussian-process
        surrogates~\cite{rasmussen_gp_2006}, executes
        NSGA-II~\cite{deb_nsga2_2002} on them and selects the infill
        batch. Steps 2 and 3 repeat for each infill iteration, and the
        checkpoint stores every evaluation.
  \item The post-analysis re-evaluates the designs of the checkpoint
        outside the loop, and its eight-layer cycle lengths define the
        constraint of the continuation of the search, which returns to
        steps 2 and 3 with its own checkpoint.
\end{enumerate}

\section{Blocks}\label{sec:arch-blocks}

\subsection{Environment and nuclear data}

Every physics script is executed inside one conda environment,
\texttt{openmc-env}, fixed by \file{environment.yml} and reproduced in a
Docker image by the \file{Dockerfile}. The environment provides OpenMC
0.15.3~\cite{romano_openmc_2015,romano_2021_depletion}, pymoo
0.6.2~\cite{blank_deb_2020} and scikit-learn~\cite{pedregosa_sklearn_2011}.
The nuclear data are the ENDF/B-VII.1 continuous-energy cross
sections~\cite{chadwick_endf71_2011} and the depletion chain
\file{chain_endfb71_pwr.xml}, downloaded once per machine by
\file{setup_data.sh}. Every script finds them through the environment
variables \texttt{OPENMC\_CROSS\_SECTIONS} and
\texttt{OPENMC\_CHAIN\_FILE}, so no file path is written in the source.

\subsection{Core model}

The five modules of Table~\ref{tab:arch-builders} build the OpenMC models
in memory from a design dictionary. They have no command line and write
no result file.

\begin{table}[htbp]
  \centering\footnotesize
  \caption{Core model: modules imported by the physics scripts.}
  \label{tab:arch-builders}
  \begin{tabularx}{\linewidth}{@{}>{\raggedright\arraybackslash}p{3.6cm} X@{}}
    \toprule
    Module & Role \\
    \midrule
    \file{reactor_model.py} & Materials, fuel pin, guide tube, control rod
      universes, $17\times17$ assembly and 32-assembly core, gadolinia
      pin pattern \\
    \file{zoning.py} & Ring map and enrichment multipliers of the zoned
      core, control bank positions, core solve at beginning of life \\
    \file{core_geometry.py} & Vessel envelope, geometric constraint,
      end-of-cycle crossing of the reactivity history, interpolation of
      the $k_\mathrm{target}$ table \\
    \file{leu_policy.py} & Enrichment limit and bounds of the design
      space, defined once for the evaluator and the optimiser \\
    \file{hardware3d.py} & Three-dimensional core with the axial stack of
      the assembly, the control rod stack, the core barrel, the
      downcomer and the vessel wall \\
    \bottomrule
  \end{tabularx}
\end{table}

\subsection{Design evaluator}

The evaluator, \file{openmc_evaluator.py}, maps a design vector to the
objectives and the constraints of Campaign 9:

\[
  \mathbf{x} = \left(e,\, w_{\mathrm{Gd}},\, t_{\mathrm{refl}},\, n_{\mathrm{Gd}}\right)
  \quad\longmapsto\quad
  \left(F_{\Delta H},\, c_\mathrm{max},\, \mathbf{g}\right)
\]

\noindent where $e$ is the enrichment, $w_{\mathrm{Gd}}$ the gadolinia
weight fraction, $t_{\mathrm{refl}}$ the reflector thickness,
$n_{\mathrm{Gd}}$ the number of poisoned pins, and $\mathbf{g}$ the
constraint vector of Section~\ref{sec:meth-c9}, each entry normalised by
its limit. For each design it executes the assembly depletion, which
gives the cycle length through the end-of-cycle crossing of
$k_\mathrm{target}$ read from \file{ktarget_table_c8.json}, and the core
solves at beginning of life, unrodded at 1000, 2000 and, where needed,
\SI{3000}{ppm}, and with the regulating banks inserted. The module
\file{boron_objective.py} turns the unrodded solves and the hump of the
depletion history into $c_\mathrm{max}$ with
Equation~\eqref{eq:c9-cmax}, and it has no OpenMC dependency, so the
post-analysis reuses it.

\subsection{Optimiser}

Table~\ref{tab:arch-optimiser} lists the scripts of the loop and of the
continuation of the search. The checkpoint of Campaign 9 is
\file{out_c9/optimization_checkpoint.json}, and that of the continuation
\file{out_c9a/optimization_checkpoint.json}.

\begin{table}[htbp]
  \centering\footnotesize
  \caption{Optimiser: scripts of the loop and of the continuation.}
  \label{tab:arch-optimiser}
  \begin{tabularx}{\linewidth}{@{}>{\raggedright\arraybackslash}p{4.3cm} X >{\raggedright\arraybackslash}p{3.4cm}@{}}
    \toprule
    Script & Role & Output \\
    \midrule
    \file{reactor_optimization.py} & Design space, Gaussian-process
      surrogates, NSGA-II on the surrogates, acquisition function,
      hypervolume & \file{optimization_}\newline\file{checkpoint.json} \\
    \file{run_optimization.py} & Command-line program: design of
      experiments, infill iterations, resume, stored settings & execution log \\
    \file{axial_ratio_model.py} & Fit of the ratio of the eight-layer to
      the assembly cycle length, Equation~\eqref{eq:axial-floor} &
      minimum assembly cycle length per design \\
    \file{seed_c9_floor.py} & Campaign 9 designs rescored under the
      constant minimum of stage 1 & seed checkpoint \\
    \file{seed_c9_axial.py} & Designs rescored under the rule of
      stage 2 & seed checkpoint \\
    \bottomrule
  \end{tabularx}
\end{table}

\subsection{Post-analysis with transport}

The scripts of Table~\ref{tab:arch-ptr} read a checkpoint, select
designs by index or by front membership, execute OpenMC at the fidelity
given on their command line, store every solve in a \file{runs.json}
keyed by design and configuration, and write a \file{summary.json}. The
store makes every script resumable.

\begin{table}[htbp]
  \centering\footnotesize
  \caption{Post-analysis of Campaign 9: scripts that execute OpenMC.}
  \label{tab:arch-ptr}
  \begin{tabularx}{\linewidth}{@{}>{\raggedright\arraybackslash}p{4.3cm} X >{\raggedright\arraybackslash}p{1.6cm}@{}}
    \toprule
    Script & Role & Section \\
    \midrule
    \file{dep_common.py} & Functions shared by the three depletion scripts
      below & -- \\
    \file{c9_dep_replicas.py} & Seed replicates and high-fidelity
      assembly depletion of the front &
      \ref{sec:res-c9-dep-check} \\
    \file{c9_dep_core2d.py} & Depletion of the zoned two-dimensional core
      & \ref{sec:res-c9-axial} \\
    \file{c9_dep_core3d.py} & Depletion of the three-dimensional core in
      $N$ axial layers, for the 8 designs, the continuation, the
      rod-count study and the seed replicates & \ref{sec:res-c9-gd},
      \ref{sec:res-c9-axial} \\
    \file{axial_shape_c9.py} & Axial power shape, burnup profile and
      peaking factors over the cycle & \ref{sec:res-c9-axial} \\
    \file{confirm3d.py} & Three-dimensional core with the structural
      components in 3 rod configurations, 2 seeds &
      \ref{sec:res-c9-ctrl3d}, \ref{sec:res-c9-peaking} \\
    \file{ladder_core_c9.py} & Fidelity study of the core peaking
      estimator & \ref{sec:res-c9-peaking} \\
    \file{mtc_scan.py} & Moderator temperature coefficient against boron
      at 2 temperatures & \ref{sec:res-c9-ceiling} \\
    \file{boron_worth.py} & Boron worth over the scanned concentrations &
      \ref{sec:res-c9-front} \\
    \file{c9_eval_predicted_front.py} & Transport evaluation of the front
      predicted by the surrogates & \ref{sec:res-c9-validation} \\
    \bottomrule
  \end{tabularx}
\end{table}

\subsection{Post-analysis without transport}

The scripts of Table~\ref{tab:arch-pnt} read the checkpoints and the
stored results of the transport scripts, and execute no transport
calculation.

\begin{table}[htbp]
  \centering\footnotesize
  \caption{Post-analysis of Campaign 9: scripts that read stored results.}
  \label{tab:arch-pnt}
  \begin{tabularx}{\linewidth}{@{}>{\raggedright\arraybackslash}p{4.3cm} X >{\raggedright\arraybackslash}p{1.6cm}@{}}
    \toprule
    Script & Role & Section \\
    \midrule
    \file{c9_front.py} & Feasible designs, front, two-bank designs &
      \ref{sec:res-c9-front} \\
    \file{c9_surrogate_cv.py} & Leave-one-out and walk-forward validation
      of the surrogates & \ref{sec:res-c9-validation} \\
    \file{c9_step0.py} & Front membership under the transport noise &
      \ref{sec:res-c9-validation} \\
    \file{c9_gp_predicted_fronts.py} & Front predicted by the surrogates
      after each refit & \ref{sec:res-c9-validation} \\
    \file{c9_acquisition_replay.py} & A posteriori re-execution of the
      acquisition under 6 selection rules & \ref{sec:res-c9-acq} \\
    \file{c9_replay_fidelity.py} & Agreement of the re-execution with the
      evaluated designs & \ref{sec:res-c9-acq} \\
    \file{c9_reverse_retro.py} & Campaign 9 designs under the objectives of
      Campaign 8 & \ref{sec:res-c9-compare} \\
    \file{c9_peaking_2d3d.py} & Two- against three-dimensional peaking
      factor & \ref{sec:res-c9-peaking} \\
    \file{c9_axial_checks.py} & Statistical checks of the layered
      depletions & \ref{sec:res-c9-axial} \\
    \file{c9_axial_correction.py} & Gadolinia-dependent cycle-length
      ratio, Equation~\eqref{eq:axial-gd} & \ref{sec:res-c9-axial} \\
    \file{c9_axial_rescore.py} & Front before and after the axial
      correction & \ref{sec:res-c9-axial} \\
    \file{mtc_front_table.py} & Three-point fit of each MTC scan, limit and
      margin of each lattice & \ref{sec:res-c9-ceiling} \\
    \file{c9_repro_check.py} & Reproducibility of the depletion with the
      campaign seed & \ref{sec:meth-fid-seed} \\
    \bottomrule
  \end{tabularx}
\end{table}

\section{Interfaces between blocks}\label{sec:arch-interfaces}

Table~\ref{tab:arch-interfaces} lists the files that cross the block
boundaries of Figure~\ref{fig:code-architecture}. Every file is plain
text, JSON or CSV, so any stage can be inspected without the code.

\begin{table}[htbp]
  \centering\footnotesize
  \caption{Code of Campaign 9: files exchanged between the blocks.}
  \label{tab:arch-interfaces}
  \begin{tabularx}{\linewidth}{@{}>{\raggedright\arraybackslash}p{4.6cm} >{\raggedright\arraybackslash}p{2.4cm} >{\raggedright\arraybackslash}p{2.4cm} X@{}}
    \toprule
    File & Producer & Consumer & Content \\
    \midrule
    \file{ktarget_table_c8.json} & Campaign 8 & Evaluator &
      $k_\mathrm{target}$ against reflector thickness, with the axial
      factor $L_{\mathrm{ax}}$ \\
    design vector $\mathbf{x}$ & Optimiser & Evaluator & One design per
      call \\
    $(F_{\Delta H}, c_\mathrm{max}, \mathbf{g})$ & Evaluator & Optimiser &
      Objectives and normalised constraints \\
    \file{out_c9/optimization_}\newline\file{checkpoint.json} & Optimiser &
      Post-analysis & Every evaluation of Campaign 9 with its
      settings \\
    \file{out_c9a/optimization_}\newline\file{checkpoint.json} & Optimiser &
      Post-analysis & Every evaluation of the continuation \\
    \file{runs.json}, \file{summary.json} & Post-analysis with transport &
      Post-analysis without transport & Solves and results of each study \\
    eight-layer cycle lengths & Post-analysis & Optimiser & Constraint of the
      continuation of the search \\
    \bottomrule
  \end{tabularx}
\end{table}

The checkpoint is the central file. It is written after every infill
iteration and stores the design space, the constraint definition, the
transport fidelity, the $k_\mathrm{target}$ table and the command line.
A resumed execution compares these fields with its own arguments and
warns on any difference. Every post-analysis script reads the
checkpoint and not the case directories, so a campaign can be analysed
again from a single file.
